\section{Claude 模型家族与工程实践}

\subsection{命名体系与迭代策略}

以诗歌形式命名，对应不同的性能-成本权衡：
\begin{itemize}
    \item \textbf{Haiku}：最小、最快、最便宜——\textit{``出奇地聪明''}
    \item \textbf{Sonnet}：中等——更聪明但更慢更贵
    \item \textbf{Opus}：最大、最强——发布时最聪明的版本
\end{itemize}

核心策略是\textbf{移动权衡曲线}：每一代新模型不只是在某个点上更好，而是整条曲线向右上方移动。Sonnet 3.5 已经比原始 Opus 3 更聪明（尤其在代码方面），Haiku 3.5 大致等于 Opus 3。\textit{``目标是移动那条曲线。''}

命名本身也是一个出人意料的挑战。AI 模型不像传统软件（3.7, 3.8 递增版本），因为预训练改进可能比预期快，打乱了命名计划。\textit{``没有人搞清楚命名……这不知怎么是一个和普通软件不同的范式。''}

\subsection{模型发布间隔中的工作}

每次发布之间有大量不可见的工作：

\begin{enumerate}
    \item \textbf{预训练}：数月，使用数万张 GPU/TPU
    \item \textbf{后训练}：RLHF 及其他 RL 方法，\textit{``规模越来越大''}
    \item \textbf{安全测试}：内部 RSP 测试 + 外部机构（美国/英国 AI 安全研究所、第三方 CBRN 测试）
    \item \textbf{推理优化}：让模型在生产环境中高效运行
    \item \textbf{API 上线}：部署和上线
\end{enumerate}

\textit{``你会惊讶于多少挑战归结为软件工程和性能工程……几乎总是归结于细节，而且往往是非常非常无聊的细节。''} 不是灵光一现的 Eureka 时刻。

\subsection{``Claude 变笨了吗？''——Wi-Fi 心理效应}

这是 Reddit 上最常见的抱怨之一。Dario 明确澄清：在没有宣布新模型的情况下，\textbf{模型权重不会改变}。

\begin{warningbox}{为什么所有人都觉得模型变笨了}
这是一个\textbf{普遍现象}——所有前沿模型公司（GPT-4、Claude、Gemini）都面临同样的投诉。可能的解释：

\begin{itemize}
    \item \textbf{心理效应}：新鲜感消退后，缺点更显眼
    \item \textbf{措辞敏感}：微小的表述变化可能导致不同输出
    \item \textbf{System Prompt 变更}：Artifacts 功能默认开启改变了 System Prompt
    \item \textbf{选择性记忆}：记住不好的体验，忘记好的
\end{itemize}

Dario 用飞机 Wi-Fi 做了精彩类比：\textit{``第一次在飞机上有 Wi-Fi 时，感觉像魔法；现在则是'这破东西怎么又不行了，真是垃圾'。''} 基线提高后，期望也随之提高。

AB 测试只在模型发布前后短暂运行。Dario 承认 System Prompt 偶尔会变化，但权重本身是不变的。
\end{warningbox}

\subsection{Claude 性格中的``打地鼠''难题}

管理 Claude 的行为是一个多维优化问题：修复一个问题往往引发另一个。

\begin{knowledgebox}{行为控制的内在困难}
修复冗长 $\to$ 模型变懒（代码中写``其余代码在这里''）。减少拒绝 $\to$ 安全问题增多。

\textit{``不是因为我们想节省算力……而是控制模型行为真的很难。''} 这是一个\textit{``未来 AI 控制问题在当下的类比''} ——如果我们今天就难以让模型在``帮你上研究生病毒学课''和``不帮你制造天花''之间划线，未来控制更强大的 AI 会有多难？

\textit{``你让模型不帮人制造和传播天花，但它愿意在你的研究生级病毒学课上帮你——如何同时做到这两点？这很难。''}
\end{knowledgebox}

\subsection{本章小结}

Claude 家族通过不断移动性能-成本权衡曲线来迭代。``变笨''是心理效应而非真实退化。模型行为的``打地鼠''难题是 AI 控制问题的当下缩影。

